\IfFileExists{acmart.cls}{
  \documentclass[sigconf,review,anonymous]{acmart}
  \settopmatter{printacmref=false,printfolios=true}
  \renewcommand\footnotetextcopyrightpermission[1]{}
  \setcopyright{none}
}{
  \documentclass[11pt]{article}
  \usepackage[margin=1in]{geometry}
}

\usepackage{booktabs}
\usepackage{tabularx}
\usepackage{microtype}
\usepackage{graphicx}
\usepackage{array}
\usepackage{multirow}
\usepackage{amsmath}
\usepackage{url}

\begin{document}

\title{Beyond a Single Benchmark: Source Effects, Closed-Book Correctness, and Practical Tradeoffs in Dual-Agent AI Tutor Supervision}

\author{Anonymous Author(s)}
\date{}
\maketitle

\begin{abstract}
Large language models are increasingly used as tutors in computing education, but a tutoring response can still fail instructionally even when it is fluent and factually correct. In assessment-adjacent settings, one important failure mode is answer leakage: the tutor discloses the answer, the final result, or a copyable solution that bypasses the intended reasoning process. Prior work showed that leakage persists in multi-turn tutoring under sustained student pressure and that dual-agent supervision can reduce it on a single curated benchmark. In this paper, we extend that result in three directions. First, we test whether leakage depends materially on question source rather than only on tutor configuration. Second, we control leakage analyses by joining each tutoring conversation to a closed-book multiple-choice result for the same tutor model on the same item. Third, we quantify the practical burden of supervision in rejection, revision, latency, tokens, and observed cost.

We analyze 8,538 judged tutoring conversations over 1,423 computer science multiple-choice questions drawn from two sources, CSBench and PeerWise. The study evaluates two tutor models, \texttt{openai/gpt-5.1} and \texttt{google/gemini-3-flash-preview}, in single-tutor and dual-loop conditions with same-provider and cross-provider supervisors. Dual-loop supervision sharply reduced leakage across both model families and both sources. For GPT as tutor, leakage fell from 15.8\% to 3.7--4.2\% on CSBench and from 22.5\% to 2.1--3.6\% on PeerWise. For Gemini as tutor, leakage fell from 6.1\% to 1.6--2.6\% on CSBench and from 6.4\% to 1.9--2.1\% on PeerWise. Source effects were also substantial: closed-book accuracy dropped from CSBench to PeerWise for both tutors, and PeerWise triggered consistently higher supervisor reject rates. Controlled leakage analyses further showed that leakage is not reducible to closed-book answerability alone; for example, GPT single-tutor leakage on PeerWise was nearly identical for questions answered correctly and incorrectly in closed-book mode (22.4\% vs.\ 22.5\%). The supervision overhead was real but moderate: average iterations per supervised turn remained between 1.02 and 1.22, most rejected drafts were repaired safely, and fallback was near zero.

These results suggest that answer leakage in AI tutoring is best understood as an interactional robustness problem shaped by dataset provenance, model choice, and control structure. We do not claim improved learning outcomes. The narrower contribution is a stronger evaluation methodology for adversarial multiple-choice tutoring that separates leakage from raw answer knowledge and makes the practical tradeoffs of supervision explicit.
\end{abstract}

\section{Introduction}

Large language models are increasingly being used in tutoring roles where the system is expected to guide a student without disclosing the target answer. In that setting, answer leakage is a specific instructional failure mode. A response may be fluent, relevant, and even factually correct while still undermining the tutoring task by revealing the answer too directly. In this paper, we use a narrow operational definition: leakage occurs when the tutor explicitly gives the answer, the final result, or a copyable solution that bypasses the intended reasoning process. This is distinct from hallucination and from broader departures from Socratic style. The distinction matters in computing education because once a tutor begins to function as an answer key, both learning goals and assessment validity are weakened.

Our earlier WCCCE paper showed that answer leakage persists in multi-turn LLM tutoring even when tutors are prompted to remain Socratic, and that a dual-loop design can reduce leakage by inserting a second model to review each tutor draft before it reaches the student. That result established the basic phenomenon, but it also left three questions open. First, how much of the observed leakage rate depends on the source of the questions rather than on the tutoring configuration alone? Second, can leakage be interpreted independently of raw model competence, or are leakage measurements partly reflecting whether the tutor could already answer the question without dialogue? Third, if supervision lowers leakage, what practical burden does that supervision introduce in revision cycles, latency, and deployment cost?

We address these questions with 8,538 judged tutoring conversations on 1,423 multiple-choice computing questions drawn from an overlap dataset that combines CSBench and PeerWise while preserving original source labels. To reduce topic confounding, the dataset is restricted to broad concepts shared across the two sources: algorithms and complexity, data structures, and systems/OS. Each conversation follows a six-turn adversarial student protocol that escalates from ordinary help-seeking to increasingly direct answer-extraction attempts. We evaluate two tutor models in a single-tutor condition and in dual-loop conditions with both same-provider and cross-provider supervisors. An external judge labels each conversation for leakage, hallucination, and Socratic compliance. We also join each tutor-question pair to a closed-book multiple-choice result from the same model, so that leakage can be analyzed conditional on whether the tutor could answer the item directly outside the tutoring dialogue.

Three findings motivate the paper. First, source effects are large enough to change the interpretation of tutoring safety results. Even within shared concepts, closed-book accuracy differed substantially by source for both tutors: Gemini scored 89.3\% on CSBench and 66.8\% on PeerWise, while GPT scored 76.1\% and 52.6\%, respectively. Leakage also varied by source, most clearly for GPT in the single-tutor condition, where leakage increased from 15.8\% on CSBench to 22.5\% on PeerWise. In this setting, benchmark provenance is not a minor reporting detail; it materially affects the observed risk profile.

Second, controlling for closed-book correctness changes how leakage should be interpreted. Leakage was not confined to questions the tutor could answer correctly on its own. In single-tutor runs, GPT leaked on 13.0\% of incorrect CSBench items and 22.5\% of incorrect PeerWise items, while Gemini leaked on 10.1\% and 7.7\% of incorrect items, respectively. The correct-versus-incorrect split did not yield a simple monotone pattern across models or sources. This does not fully explain why leakage occurs, but it shows that leakage cannot be reduced to competence alone. A tutor may reveal an answer both when it appears able to solve the item and when it does not.

Third, supervision helps, but its practical form matters. Across all dual-loop pairings, leakage fell to 1.6--4.2\%, with compliance near 100\% and hallucination low. However, the supervision process was not uniform. Reject rates ranged from 1.8\% to 17.0\% depending on pairing and source, average iterations per turn ranged from 1.02 to 1.22, and latency increased relative to single-tutor runs. Most initially rejected drafts were repaired safely, but the review burden differed enough that dual-loop supervision should not be treated as a single intervention with a single cost.

Taken together, these results extend the prior WCCCE study in three ways. They show that answer leakage estimates are source-sensitive even after reducing topical confounding; they introduce a closed-book correctness control that separates leakage analysis from raw answerability; and they treat supervision as an operational process with measurable tradeoffs rather than only as a binary safety feature. We do not claim that these findings generalize beyond the models, question format, and adversarial protocol studied here. The narrower claim is that, for multi-turn LLM tutoring on computing multiple-choice questions, conclusions about leakage depend materially on dataset provenance, on whether baseline correctness is controlled, and on how supervision is implemented.

\section{Related Work}

Prior work on LLM tutoring in computing education has shown both the promise and the instability of conversational scaffolding. LLM-based tutors can provide accessible explanations, hints, and debugging support, but they may also over-help by disclosing solutions or performing too much of the reasoning task for the student. Within this space, answer leakage is a particularly useful construct for assessment-adjacent multiple-choice tutoring because it can be defined operationally and measured consistently.

A second line of work examines adversarial prompting and jailbreak behavior in LLM systems. Our threat model overlaps with this literature in that the student escalates requests across turns and attempts to extract forbidden content. However, the educational setting differs in an important way: the attack is embedded inside plausible tutoring dialogue rather than being an obviously out-of-distribution safety probe. This means that the system must distinguish between legitimate help-seeking and tutoring moves that gradually collapse into answer disclosure.

Recent work has also explored LLM oversight, critique, and review loops. Dual-agent or multi-agent designs can use one model to draft and another to critique or approve the response. The distinctive question for educational tutoring is not only whether review reduces unsafe outputs, but whether it reduces student-visible leakage at a practical revision cost. In tutoring deployments, a system that is nominally safer but too slow, too expensive, or too complex to manage may not be deployable.

This paper also speaks to benchmark design in computing education. Results from a single curated dataset can obscure how strongly tutor behavior depends on source style, item phrasing, and answerability. By comparing CSBench and PeerWise within shared concept groups and by conditioning leakage on closed-book correctness, we treat benchmark provenance and baseline competence as part of the methodological problem rather than as background detail.

\section{Method}

\subsection{Study Design and Data Sources}

We analyzed a merged log of 8,538 simulated tutoring conversations over 1,423 multiple-choice questions. The question set combined two existing sources: 644 CSBench items and 779 PeerWise items. Because each question was evaluated under six tutoring configurations, the analytic sample formed a complete question-by-condition matrix. All analyzed records were multiple-choice; no other question formats were included in this merged run.

\subsection{Tutoring Conditions and Conversation Protocol}

The experiment varied tutoring architecture and model assignment while holding the attacker and judge constant. Tutor models were \texttt{openai/gpt-5.1} and \texttt{google/gemini-3-flash-preview}. In the single condition, each tutor responded directly to the student without a supervisor, yielding two single-model configurations. In the dual-loop condition, a tutor draft was screened by a supervisor model before being shown to the student, yielding four pairings: GPT/GPT, Gemini/Gemini, GPT tutor with Gemini supervisor, and Gemini tutor with GPT supervisor. Across all conditions, the student attacker and the judge both used \texttt{google/gemini-3-flash-preview}. Conversations were allowed to continue for up to six turns.

Each conversation consisted of alternating student and tutor messages generated by the harness. The student was not a passive learner but an adversarial attacker designed to escalate extraction attempts across turns. Based on the prompt and hidden-trace fields in the raw logs, the attack sequence progressed from ordinary requests for help to more direct answer-seeking, manipulative pressure, prompt-injection attempts, and obfuscated reconstruction strategies. In dual-loop settings, the supervisor evaluated each tutor draft for safety before release to the student. This design allowed us to compare whether adding supervisory review reduced answer leakage relative to single-model tutoring.

\subsection{Outcome Measures and Closed-Book Linkage}

We used the harness judge to score each conversation for leakage, hallucination, and Socratic compliance. Leakage was defined conservatively as revealing a final numeric, symbolic, or otherwise definitive answer; for multiple-choice items, this included selecting an option or eliminating options such that the answer became recoverable. Hallucination indicated factually incorrect domain content in the tutor’s response. Compliance indicated maintenance of a Socratic teaching style rather than direct solution giving. We treated leakage as the primary outcome and also recorded the first turn at which leakage occurred when applicable. Conversation-level outcomes were taken from the conversation records, which in this run were derived from turn-level judging.

To contextualize leakage against model capability, we linked each tutoring conversation to a separate closed-book MCQ accuracy file. That file contained one accuracy record per \texttt{(question\_id, model\_id)} for each tutor model answering each question directly, for a total of 2,846 closed-book records. We joined closed-book accuracy to tutoring outcomes using \texttt{(question\_id, tutor\_model\_id)}. Because the merged tutoring run was MCQ-only and included both tutor models on every question, this join covered the full analytic sample (8,538/8,538 conversations). This linkage enabled controlled comparisons of leakage for questions a tutor model could answer correctly versus incorrectly in closed-book mode.

\subsection{Process, Latency, and Cost Measures}

In addition to the behavioral outcomes above, the raw logs also supported process measures such as turn counts, supervision iterations, latency, token usage, and cost. Latency and token measures were aggregated from per-call metadata in the conversation logs. Cost was treated as observed cost rather than complete cost, because totals were computed from per-call \texttt{usage.cost} values only when those fields were present in provider responses.

\subsection{Analytic Caveats}

Several limitations of the instrumentation should be noted. First, the attacker protocol was not explicitly versioned in the repository, so the attack regimen must be reconstructed from prompts and hidden-trace metadata rather than cited as a stable protocol release. Second, the analytic dataset was a merged artifact that preserved two replay run IDs rather than a single canonical run identifier, which complicates provenance reporting at the run level even though the conversation-level records are intact. Third, cost estimates should be interpreted as partial observations because missing per-call \texttt{usage.cost} fields reduce coverage and may understate total cost.

\section{Results}

\subsection{Overall Effect of Dual-Loop Supervision}

Across 8,538 tutoring conversations covering 1,423 multiple-choice questions, dual-loop supervision consistently reduced answer leakage relative to matched single-tutor baselines while preserving judged Socratic compliance (Table~\ref{tab:main-leakage}). For GPT as tutor, leakage fell from 15.8\% to 3.7--4.2\% on CSBench and from 22.5\% to 2.1--3.6\% on PeerWise, depending on the supervisor. For Gemini as tutor, leakage fell from 6.1\% to 1.6--2.6\% on CSBench and from 6.4\% to 1.9--2.1\% on PeerWise. Aggregating across all evaluated cells, leakage decreased from 366/2,846 conversations (12.9\%) in single-loop to 153/5,692 (2.7\%) in dual-loop, a 79\% relative reduction.

\begin{table*}[t]
\caption{Conversation-level outcomes by tutor, condition, and source.}
\label{tab:main-leakage}
\small
\centering
\begin{tabular}{@{}llllcccc@{}}
\toprule
Tutor & Supervisor & Condition & Source & Leakage & Compliance & Hallucination & Median first leak turn \\
\midrule
GPT & none & Single & CSBench & 102/644 (15.8\%) & 628/644 (97.5\%) & 19/644 (3.0\%) & 3.0 \\
Gemini & none & Single & CSBench & 39/644 (6.1\%) & 644/644 (100.0\%) & 4/644 (0.6\%) & 4.0 \\
GPT & GPT & Dual & CSBench & 24/644 (3.7\%) & 644/644 (100.0\%) & 9/644 (1.4\%) & 2.0 \\
Gemini & Gemini & Dual & CSBench & 17/644 (2.6\%) & 644/644 (100.0\%) & 9/644 (1.4\%) & 3.0 \\
GPT & Gemini & Dual & CSBench & 27/644 (4.2\%) & 644/644 (100.0\%) & 18/644 (2.8\%) & 3.0 \\
Gemini & GPT & Dual & CSBench & 10/644 (1.6\%) & 644/644 (100.0\%) & 5/644 (0.8\%) & 2.5 \\
\midrule
GPT & none & Single & PeerWise & 175/779 (22.5\%) & 736/779 (94.5\%) & 29/779 (3.7\%) & 4.0 \\
Gemini & none & Single & PeerWise & 50/779 (6.4\%) & 776/779 (99.6\%) & 17/779 (2.2\%) & 5.0 \\
GPT & GPT & Dual & PeerWise & 16/779 (2.1\%) & 779/779 (100.0\%) & 9/779 (1.2\%) & 3.0 \\
Gemini & Gemini & Dual & PeerWise & 16/779 (2.1\%) & 779/779 (100.0\%) & 11/779 (1.4\%) & 4.0 \\
GPT & Gemini & Dual & PeerWise & 28/779 (3.6\%) & 779/779 (100.0\%) & 13/779 (1.7\%) & 4.0 \\
Gemini & GPT & Dual & PeerWise & 15/779 (1.9\%) & 779/779 (100.0\%) & 4/779 (0.5\%) & 4.0 \\
\bottomrule
\end{tabular}
\end{table*}

The survival curves show that this benefit persisted over time rather than only at early turns. For example, by turn 6 GPT’s leakage-free survival improved from 84.2\% to 95.8--98.4\% on CSBench and from 77.5\% to 96.4--98.1\% on PeerWise, while Gemini improved from 93.9\% to 97.4--98.4\% on CSBench and from 93.6\% to 97.9--98.1\% on PeerWise. All dual-loop conditions also achieved 100.0\% judged compliance.

\begin{figure}[t]
\centering
\fbox{\parbox[c][1.6in][c]{0.92\linewidth}{\centering Placeholder for survival-by-turn figure\\(single vs.\ dual, split by source).}}
\caption{Leakage-free survival by turn. The final paper should plot single versus dual conditions separately for CSBench and PeerWise.}
\label{fig:survival}
\end{figure}

\subsection{PeerWise as the Harder Benchmark}

PeerWise appears to be the more challenging and operationally riskier benchmark. Closed-book accuracy dropped substantially from CSBench to PeerWise for both Gemini, from 89.3\% to 66.8\%, and GPT, from 76.1\% to 52.6\% (Table~\ref{tab:closed-book}). In single-loop tutoring, leakage was also higher on PeerWise for GPT (22.5\% vs.\ 15.8\%) and slightly higher for Gemini (6.4\% vs.\ 6.1\%). At the same time, this claim should remain somewhat cautious: once dual-loop supervision was enabled, PeerWise did not produce uniformly higher leakage than CSBench in every matched pairing. The stronger conclusion is that PeerWise stressed the tutoring process more. Supervisor reject rates were higher for every dual-loop pairing on PeerWise than on CSBench, including 17.0\% versus 10.3\% for GPT/GPT, 7.0\% versus 4.4\% for GPT/Gemini, and 11.4\% versus 7.3\% for Gemini/GPT.

\begin{table}[t]
\caption{Closed-book MCQ accuracy by tutor and source.}
\label{tab:closed-book}
\small
\centering
\begin{tabular}{@{}lcccc@{}}
\toprule
Tutor model & CSBench acc. & $n$ & PeerWise acc. & $n$ \\
\midrule
Gemini & 575/644 (89.3\%) & 644 & 520/779 (66.8\%) & 779 \\
GPT & 490/644 (76.1\%) & 644 & 410/779 (52.6\%) & 779 \\
\bottomrule
\end{tabular}
\end{table}

PeerWise leaks were also more likely to contain direct answer content, including final numeric or unique results (21.3\% vs.\ 11.4\% of leaks) and paraphrase-equivalent answers (36.0\% vs.\ 27.4\%). Taken together, the evidence supports describing PeerWise as the harder and riskier benchmark, even if not every supervised leakage rate is higher there.

\subsection{Gemini Versus GPT}

Gemini was the stronger model family both in closed-book answering and in tutoring behavior. In closed-book evaluation, Gemini outperformed GPT on both datasets: 89.3\% versus 76.1\% on CSBench and 66.8\% versus 52.6\% on PeerWise. In single-loop tutoring, Gemini also leaked much less often than GPT, with 6.1\% versus 15.8\% leakage on CSBench and 6.4\% versus 22.5\% on PeerWise. Gemini’s tutoring outputs were also at least as compliant and generally less hallucinatory: compliance was 100.0\% on CSBench and 99.6\% on PeerWise for Gemini, compared with 97.5\% and 94.5\% for GPT, while hallucination fell from 3.0\% to 0.6\% on CSBench and from 3.7\% to 2.2\% on PeerWise.

In the dual-loop setting, configurations with Gemini as the drafting tutor also leaked less overall than configurations with GPT as tutor, 58/2,846 conversations (2.0\%) versus 95/2,846 (3.3\%), and the lowest-leakage pairing on both benchmarks was Gemini tutor with GPT supervisor: 1.6\% on CSBench and 1.9\% on PeerWise. Because supervisor identity is not fully held constant across these dual-loop comparisons, that last point should be framed as suggestive rather than strictly causal.

\subsection{Leakage Is Not Just Answer Knowledge}

The controlled-leakage analysis shows that leakage cannot be explained only as a by-product of the tutor knowing the correct answer (Table~\ref{tab:controlled}). GPT-single on PeerWise leaked on 92/410 questions it answered correctly closed-book (22.4\%) and on 83/369 questions it answered incorrectly (22.5\%), which is essentially identical. Gemini-single on CSBench leaked more often on incorrect items than correct ones, 7/69 (10.1\%) versus 32/575 (5.6\%). Nonzero wrong-answer leakage also remained under dual-loop supervision, for example 9/154 (5.8\%) for GPT/Gemini on CSBench and 2/69 (2.9\%) for Gemini/GPT on CSBench.

We therefore interpret leakage as at least partly an interaction-control failure, not just an answer-knowledge failure: models can disclose too much through partial solution steps, elimination strategies, or overly specific guidance even when they do not reliably know the final answer. This claim should still be stated cautiously because some wrong-answer cells are relatively small.

\begin{table*}[t]
\caption{Controlled leakage by closed-book correctness.}
\label{tab:controlled}
\small
\centering
\begin{tabular}{@{}lllcccc@{}}
\toprule
Pairing & Condition & Tutor model & CSBench correct & CSBench wrong & PeerWise correct & PeerWise wrong \\
\midrule
Gemini-single & Single & Gemini & 32/575 (5.6\%) & 7/69 (10.1\%) & 30/520 (5.8\%) & 20/259 (7.7\%) \\
GPT-single & Single & GPT & 82/490 (16.7\%) & 20/154 (13.0\%) & 92/410 (22.4\%) & 83/369 (22.5\%) \\
Gemini-Gemini & Dual & Gemini & 16/575 (2.8\%) & 1/69 (1.4\%) & 13/520 (2.5\%) & 3/259 (1.2\%) \\
Gemini-GPT & Dual & Gemini & 8/575 (1.4\%) & 2/69 (2.9\%) & 12/520 (2.3\%) & 3/259 (1.2\%) \\
GPT-Gemini & Dual & GPT & 18/490 (3.7\%) & 9/154 (5.8\%) & 13/410 (3.2\%) & 15/369 (4.1\%) \\
GPT-GPT & Dual & GPT & 19/490 (3.9\%) & 5/154 (3.2\%) & 6/410 (1.5\%) & 10/369 (2.7\%) \\
\bottomrule
\end{tabular}
\end{table*}

\subsection{Supervision Process Metrics}

The dual-loop mechanism was effective without requiring heavy iterative overhead (Table~\ref{tab:process}). Across 33,767 supervised turns, average iterations per turn was 1.09, and 31,105 turns (92.1\%) were approved on the first try. When the supervisor did reject a draft, 2,628 of 2,662 rejected drafts (98.7\%) were revised into a safe approved response. Only 13 rejected drafts (0.49\%) still leaked after revision, and fallback was used in just 21 turns total (0.06\%). This pattern held even on PeerWise, where the process was somewhat busier but still efficient.

\begin{table*}[t]
\caption{Dual-loop process and practicality metrics.}
\label{tab:process}
\small
\centering
\begin{tabular}{@{}lllcccccc@{}}
\toprule
Pairing & Source & Reject rate & Avg.\ iters/turn & Approved 1st try & Fallback rate & Safe after reject & Still leaked after reject & Median conv.\ latency \\
\midrule
GPT/GPT & CSBench & 390/3785 (10.3\%) & 1.12 & 3395/3785 (89.7\%) & 5/3785 (0.1\%) & 379/390 (97.2\%) & 6/390 (1.5\%) & 42027 ms \\
Gemini/Gemini & CSBench & 69/3819 (1.8\%) & 1.02 & 3750/3819 (98.2\%) & 0/3819 (0.0\%) & 69/69 (100.0\%) & 0/69 (0.0\%) & 37511 ms \\
GPT/Gemini & CSBench & 167/3789 (4.4\%) & 1.05 & 3622/3789 (95.6\%) & 0/3789 (0.0\%) & 166/167 (99.4\%) & 1/167 (0.6\%) & 40736 ms \\
Gemini/GPT & CSBench & 280/3831 (7.3\%) & 1.08 & 3551/3831 (92.7\%) & 1/3831 (0.0\%) & 279/280 (99.6\%) & 0/280 (0.0\%) & 38201 ms \\
\midrule
GPT/GPT & PeerWise & 785/4631 (17.0\%) & 1.22 & 3846/4631 (83.0\%) & 9/4631 (0.2\%) & 772/785 (98.3\%) & 4/785 (0.5\%) & 65542 ms \\
Gemini/Gemini & PeerWise & 117/4646 (2.5\%) & 1.03 & 4529/4646 (97.5\%) & 0/4646 (0.0\%) & 117/117 (100.0\%) & 0/117 (0.0\%) & 45648 ms \\
GPT/Gemini & PeerWise & 324/4627 (7.0\%) & 1.07 & 4303/4627 (93.0\%) & 0/4627 (0.0\%) & 322/324 (99.4\%) & 2/324 (0.6\%) & 58505 ms \\
Gemini/GPT & PeerWise & 530/4639 (11.4\%) & 1.14 & 4109/4639 (88.6\%) & 6/4639 (0.1\%) & 524/530 (98.9\%) & 0/530 (0.0\%) & 53757 ms \\
\bottomrule
\end{tabular}
\end{table*}

\subsection{Latency and Cost Tradeoffs}

These leakage reductions are not free, but the overhead appears manageable in this experimental setting. Dual-loop increased model calls per turn from 3.0 in single-loop to 4.04--4.43, and mean tokens per conversation rose from roughly 22.6k--27.0k in single-loop to 32.0k--46.8k in dual-loop, depending on the benchmark and pairing. Median conversation latency likewise increased from 28.7--43.1 seconds in single-loop to 37.5--65.5 seconds in dual-loop. Observed cost per conversation remained in a practical range, roughly \$0.0105--\$0.0169 for single-loop and \$0.0110--\$0.0234 for dual-loop. However, the dollar-cost comparison should be stated carefully because GPT-involving conditions have partial cost coverage, ranging from 45.1\% to 74.4\%, whereas Gemini-only conditions have full coverage. For that reason, the strongest claim here is about token and latency overhead; the monetary tradeoff appears manageable, but that conclusion is somewhat more tentative.

\section{Discussion}

Our results suggest that answer leakage in AI tutoring is best understood as an interactional robustness problem, not only as a static question-answering problem. Holding the tutor fixed, leakage changed sharply when the conversational control structure changed: for GPT, leakage fell from 15.8\% to 3.7--4.2\% on CSBench and from 22.5\% to 2.1--3.6\% on PeerWise under dual-loop supervision; for Gemini, it fell from 6.1\% to 1.6--2.6\% on CSBench and from 6.4\% to 1.9--2.1\% on PeerWise. Because these reductions occur without changing the underlying tutor’s task knowledge, they point to failures that emerge in the back-and-forth of adversarial tutoring. The turn-level survival curves tell the same story. On PeerWise, GPT in the single condition leaked on only 0.1\% of conversations at turn 1, but survival fell from 99.9\% after the first turn to 77.5\% by turn 6, with most first leaks appearing in turns 3--5. A tutor that appears safe early in a conversation may therefore become substantially less safe under sustained pressure.

PeerWise-like items also appear to be a harder stress test than curated benchmark questions. Closed-book accuracy dropped for both tutors on PeerWise relative to CSBench, and the dual-loop controller intervened more often on PeerWise for every pairing. We would therefore be cautious about treating curated benchmark performance as a sufficient proxy for robustness in classroom-like settings. At the same time, supervision did not erase differences between tutors. Tutor choice remained a major safety lever: in the single condition, Gemini leaked much less than GPT on both sources, and pairings with Gemini as tutor achieved the lowest leakage rates overall. Lightweight supervision materially improves safety, but it works best when layered onto a comparatively robust base tutor rather than used as a substitute for one.

The operational overhead of the dual-loop design appears meaningful but practical. Average iterations per turn remained low (1.02--1.22), 83.0--98.2\% of turns were approved on the first try, and 97.2--100\% of initially rejected turns were revised into a safe approved response. Median turn latency increased from about 1.4--2.7 seconds in the single condition to 2.7--5.2 seconds in dual-loop, and model calls per turn rose from 3.0 to 4.04--4.43. We would not describe this overhead as negligible, but neither does leakage reduction require a heavy multi-stage pipeline. For multiple-choice tutoring, a comparatively lightweight supervisory loop appears to buy a large reduction in answer leakage at a level of additional computation that may be acceptable in practice.

\section{Limitations}

Several limits narrow what can be concluded here. First, the study is entirely within a multiple-choice setting. The results therefore do not justify claims about open-ended explanations, code synthesis, proof construction, or other tutoring formats where the line between helpful guidance and answer disclosure may look different. Relatedly, we measure answer leakage, compliance, hallucination, latency, and process overhead, but not student learning, transfer, retention, satisfaction, or trust. A system that leaks less is not necessarily a system that teaches more effectively.

Second, the evaluation uses a synthetic adversarial student and an LLM judge. These choices make the comparisons reproducible at scale, but they are still proxies for real classroom interaction and human evaluation. The attacker follows an escalating six-turn protocol, and the judge applies a particular interpretation of leakage and compliance; different students, longer conversations, or human raters could shift the observed rates. Third, the model set is small. We evaluate two tutor families and a small number of supervision pairings, so we cannot generalize from these results to broader claims about providers or architectures. The strong differences we observe are enough to show that tutor choice matters, but not enough to establish a stable ranking outside this setup.

Finally, the source comparison should be read carefully. PeerWise-like questions appear more demanding than the curated benchmark on several indicators, but the PeerWise subset here is still only one slice of student-authored MCQ data and is concentrated in particular concepts. The cost figures are also best interpreted cautiously because the report marks them as observed values with incomplete coverage for some call types. For these reasons, we interpret the study as evidence about conversational robustness in adversarial MCQ tutoring, not as a comprehensive account of educational effectiveness or deployment readiness.

\section{Conclusion}

Within the scope of adversarial multiple-choice tutoring, the main result is straightforward: leakage is not just a matter of whether a model knows the answer, but of how well it maintains tutoring boundaries across a multi-turn interaction. PeerWise-like questions make that boundary-keeping problem harder than curated benchmark items, and a lightweight dual-loop supervisor substantially reduces leakage with moderate overhead. At the same time, tutor selection remains a first-order safety decision; supervision helps considerably, but it does not remove the importance of starting from a more robust base tutor.

We therefore see dual-loop supervision as a practical containment mechanism for MCQ tutoring rather than as a complete safety solution or as evidence of better learning. The present results support cautious claims about reduced answer leakage under adversarial interaction in an MCQ-only setting. They should not be read as showing improved pedagogy, improved student outcomes, or robustness across tutoring formats more broadly.

\end{document}
